\chapter{Counting Discipline}
\index{test counting}\index{test population}
\label{ch:counting-discipline}

``How many tests does CNA have?'' has no context-free answer.  The repository contains source
files, GoogleTest definitions, standalone executables, and CTest registrations; one renderer is
selected per configure; GoogleTest cases are discovered by executing the binary.  A useful
number must therefore name both its unit and its configuration.

\section{Four quantities that must not be exchanged}

At CNA commit \texttt{7a64362e}, static inspection gives the following representative counts:

\begin{center}
\small
\begin{tabular}{rp{9.1cm}}
\toprule
Count & Exact definition \\
\midrule
433 & C\texttt{++} source files in the three roots matched by the CnaTests glob, before
      configure-dependent filtering. \\
6,818 & GoogleTest macro definitions across all those test-directory sources. \\
6,670 & Definitions compiled into the default Linux \texttt{OPENGLES3}, Net-enabled,
        FFmpeg-available CnaTests source set. \\
1,785 & Static CTest registration statements made through the renderer-test helper, after
        expanding Skia's wrapper calls; not active tests in one build. \\
1,233 & C\texttt{++} files under module \texttt{examples} directories; 1,169 happen to be named
        \texttt{*\_test.cpp}, but almost none are GoogleTest sources. \\
\bottomrule
\end{tabular}
\end{center}

These figures describe different populations.  Adding them is meaningless, and selecting the
largest does not make it the most authoritative.  A test source can define many cases; a
parameterized definition can expand at runtime; one executable can become thousands of CTest
entries; one standalone executable can exist without any CTest registration.

\section{How the static counts are reproduced}

The source population is a sorted union, not a sum of unchecked directory totals:

\begin{verbatim}
{ find modules -path 'modules/*/tests/*' -name '*.cpp' -type f
  find tests -name '*.cpp' -type f; } | sort -u | wc -l
\end{verbatim}

Macro counts search only declaration starts, distinguishing \texttt{TEST}, \texttt{TEST\_F},
and \texttt{TEST\_P}.  The default-configure figure then subtracts the renderer test directories
excluded by CMake.  This is still a source definition count, not a run count.  Six
\texttt{INSTANTIATE\_TEST\_SUITE\_P} sites and platform conditionals can change runtime
expansion.

Renderer registrations require another derivation.  The tree contains 1,624 direct calls to
\texttt{cna\_register\_renderer\_test}; three are Skia wrapper bodies rather than registrations,
and 164 calls invoke those wrappers.  Hence $1624-3+164=1785$ static statements.  Nested CMake
conditions still decide which subset exists in any actual configure.

\section{Why there is no all-renderer total}

The build enters exactly one selected renderer family.  EasyGL's 293 registration statements,
Vulkan's 211, bgfx's 179, and Skia's 175 belong to mutually exclusive configurations.  Summing
them creates a number that no build can expose and no machine can run as one CTest suite.

CTest discovery adds a second barrier.  \texttt{PRE\_TEST} obtains case names by running
\texttt{CnaTests --gtest\_list\_tests}; configure-time source inspection cannot know the complete
expanded list.  Build features, conditional compilation, and parameter instantiation all affect
it.  The authoritative count for a particular build must come from that built artifact and must
name its renderer and feature set.

\section{Identity counts can be settled}

Not every disputed number is inherently unknowable.  The canonical cache value table contains
46 public renderer identities, while \texttt{modules/renderers} contains 42 implementation
families after excluding common infrastructure.  Four identities share implementations with
other public modes.  Thus ``46 identities, 42 families'' resolves the apparent conflict instead
of choosing one number and calling the other wrong.

This pattern is widely useful: define the noun before reconciling the count.  An identity is a
user-selectable contract; a family is a source implementation.  They need not be one-to-one.

\section{Counts that require execution}

Static inspection cannot establish:

\begin{itemize}
  \item the CTest count for a configured and built tree;
  \item how many parameterized GoogleTests expand or actually run;
  \item active standalone registrations after nested conditions;
  \item pass, failure, timing, or skip totals;
  \item whether any pixel assertion reached a display;
  \item image-oracle divergence counts;
  \item line, branch, or function coverage---no coverage tool is configured.
\end{itemize}

CNA documents several other kinds of ``coverage'': API-surface estimates, mappings from source
files to suites, and feature checklists.  None is measured code coverage.  A percentage without
that qualifier is more misleading than no percentage at all.

\section{A reporting grammar}

A durable result can fit in one sentence if every noun is present: ``On commit X, the Y renderer
build with features Z discovered N CTest entries; it ran R, skipped S, and failed F under host
H.''  A source audit uses a different sentence: ``At commit X, this command found N declarations
matching definition D.''

Three historical comments claimed approximately 247, 330, and 600 renderer registrations, each
using a different unstated meaning; the audited tree statically contains 1,785 statements across
exclusive configurations.  Comments rot, while reproducible definitions can be rerun.  The
discipline is therefore not to avoid numbers, but to make every number carry its derivation.
